\section{GPT-OSS expert reuse under a fixed-cache proxy}
\label{sec:oss-cache-diagnostics}

A separate full-diagnostics rerun evaluates GPT-OSS fixed positive identity
steering at strength $+1$, using the frozen policy and the same 1,395 held-out
attempts on 465 problems. It uses concurrency 25 and the frozen template date
2026-09-22. The compatible baseline is reused with output-hash provenance;
the guided generations are new. This diagnostic rerun is kept separate from
the historical minimal-diagnostics comparison and is not an independent
problem-set replication. No new target or strength is selected from these results.

Under the same offline answer-equivalence scoring contract as the main
intervention analysis, accuracy is $1086/1395=77.85\%$ for baseline and
$1088/1395=77.99\%$ for guidance. The change is $+0.14$ percentage points,
with a paired, dataset-stratified source-problem bootstrap 95\% interval
$[-2.37,+2.37]$ (5,000 resamples, seed 42). There are 94 wrong-to-right and
92 right-to-wrong changes. Mean generated tokens increase from 6,497.2 to
6,813.9 ($+4.88\%$), and token-limit hits increase from 74 to 77.
Thus the observed aggregate accuracy difference is small, but accuracy
preservation and an efficiency improvement are not established.

The routing counters permit a fixed resident-cache proxy at the four steered
layers, 15, 19, 21 and 23. Four target experts per layer would occupy 16
layer--expert slots in total. Define $H$ as the fraction of top-$k$ expert
selections covered by that fixed resident set; the proxy for nonresident
selections is $1-H$. Both pre-bias and post-bias counters are evaluated on
the guided trajectory's hidden states. The pre-bias counts are therefore
not routing measurements from the separately generated baseline.

\begin{table}[ht]
\centering\small
\caption{GPT-OSS diagnostic rerun: static-cache coverage as a percentage of
expert selections. Target and frequency caches each have four slots per layer.
Frequency pins are selected from calibration statistics only. All columns
use states visited by guided generation.}
\label{tab:oss-cache-coverage}
\begin{tabular}{lrrr}
\toprule
Layer & Target, pre-bias & Target, post-bias & Frequency, pre-bias \\
\midrule
15 & 5.19 & 31.40 & 30.40 \\
19 & 8.70 & 45.23 & 37.42 \\
21 & 2.20 & 22.58 & 41.36 \\
23 & 4.23 & 41.48 & 38.32 \\
All four & 5.08 & 35.17 & 36.87 \\
\bottomrule
\end{tabular}
\end{table}

Target-cache coverage rises from $5.08\%$ to $35.17\%$. Relative to pinning
those same targets without the bias on these states, this corresponds to
$31.70\%$ fewer nonresident selections:
$(H_{\mathrm{after}}-H_{\mathrm{before}})/(1-H_{\mathrm{before}})$.
However, a same-budget cache of calibration-frequent experts already covers
$36.87\%$ of pre-bias selections. Guided target pinning therefore has
$2.70\%$ more nonresident selections than this control. A hindsight-optimal
static cache covers $37.39\%$ before the bias; against it, guided target
pinning has $3.55\%$ more nonresident selections. The hindsight cache is an
analytical control, not a deployable policy or a bound on adaptive caching.

The supported result is greater reuse of the boosted expert identities,
alongside a small observed aggregate accuracy difference. For this OSS
policy, that concentration does not outperform the calibration-frequency
static-cache control in the measured proxy. This distinction prevents
increased target frequency alone from being interpreted as a systems gain.

These counters include prefill, decode and engine padding, cover only the
four selected layers, and compute top-$k$ from hook logits rather than
separately observing native dispatch. They retain neither temporal order
nor per-attempt routing, so no cache-effect confidence interval is available.
The proxy assumes an already resident pinned set and excludes cold-start
loads, dynamic caching, shared loads across batched tokens and prefetch.
It measures neither SSD bytes nor latency and must not be extrapolated to
whole-model speed. A small accuracy loss could be acceptable if it buys
sufficient speed, but that trade-off requires ordered decode traces and
ultimately an equal-memory offloaded-inference benchmark. Longer guided
outputs make total loads and latency per answer especially important.
